\documentclass[11pt,a4paper]{article}
\usepackage[utf8]{inputenc}
\usepackage[T1]{fontenc}
\usepackage[spanish,provide=*]{babel}
\addto\captionsspanish{\renewcommand{\tablename}{Tabla}}
\usepackage{lmodern}
\usepackage{microtype}
\usepackage[margin=2.5cm]{geometry}
\usepackage{amsmath,amssymb}
\usepackage{graphicx}
\usepackage{booktabs}
\usepackage{adjustbox}
\usepackage{pdflscape}
\usepackage{float}
\usepackage[font=small,labelfont=bf]{caption}
\usepackage{enumitem}
\usepackage{xcolor}
\usepackage{listings}
\usepackage[round]{natbib}
\usepackage[hidelinks]{hyperref}
\setlength{\parskip}{0.5em}
\setlength{\parindent}{0pt}
\setlist{itemsep=0.2em}

% Texto en rojo = algo que hay que revisar/actualizar después de volver a correr el do-file.
% Para sacar el color: \renewcommand{\revisar}[1]{#1}
\newcommand{\revisar}[1]{\textcolor{red}{#1}}
% Inserta una tabla exportada desde Stata (se achica si no entra)
\newcommand{\tabla}[1]{\begin{adjustbox}{max width=\linewidth, max totalheight=0.72\textheight}\input{#1}\end{adjustbox}}
% Figura que todavía no existe (se genera al correr el do-file): muestra un recuadro
\newcommand{\figura}[2][0.75\linewidth]{\IfFileExists{#2}{\includegraphics[width=#1]{#2}}{\fbox{\parbox{\dimexpr#1-2\fboxsep-2\fboxrule\relax}{\centering\vspace{2em}\textcolor{red}{Figura pendiente: se genera al correr el do-file}\\[0.5em]\texttt{\detokenize{#2}}\vspace{2em}}}}}
% Notas al pie de tablas
\newcommand{\notas}[1]{\par\vspace{0.3em}\begin{minipage}{0.92\linewidth}\footnotesize\textit{Notas:} #1\end{minipage}}
% Consignas en cursiva
\newcommand{\consigna}[1]{\par\textit{#1}\par}

\lstset{
  language={},
  basicstyle=\ttfamily\scriptsize,
  breaklines=true,
  columns=fullflexible,
  keepspaces=true,
  showstringspaces=false,
  frame=single,
  rulecolor=\color{gray!50},
  commentstyle=\color{gray},
  morecomment=[l]{*\ },
  morecomment=[l]{//},
  morecomment=[s]{/*}{*/},
  tabsize=4,
  inputencoding=utf8,
  extendedchars=true,
  literate={á}{{\'a}}1 {é}{{\'e}}1 {í}{{\'i}}1 {ó}{{\'o}}1 {ú}{{\'u}}1
           {Á}{{\'A}}1 {É}{{\'E}}1 {Í}{{\'I}}1 {Ó}{{\'O}}1 {Ú}{{\'U}}1
           {ñ}{{\~n}}1 {Ñ}{{\~N}}1 {ü}{{\"u}}1 {Ü}{{\"U}}1
           {¿}{{?`}}1 {¡}{{!`}}1 {’}{{'}}1 {‘}{{`}}1 {“}{{``}}1 {”}{{''}}1
           {ã}{{\~a}}1 {õ}{{\~o}}1 {â}{{\^a}}1 {ê}{{\^e}}1 {ô}{{\^o}}1 {à}{{\`a}}1
           {Ã}{{\~A}}1 {Õ}{{\~O}}1 {Â}{{\^A}}1 {Ê}{{\^E}}1 {Ô}{{\^O}}1 {Ç}{{\c C}}1 {×}{{$\times$}}1
           {°}{{$^\circ$}}1 {–}{{--}}1 {—}{{---}}1 {·}{{$\cdot$}}1 {ç}{{\c c}}1
}
\begin{document}

\begin{center}
{\Large\bfseries Problem Set 10: Regresión discontinua}\\[0.5em]
E178: Economía Aplicada --- Universidad de San Andrés\\
Fecha de entrega: 22 de noviembre de 2024\\[0.5em]
\textbf{Integrantes:} Juan Batalla, Nina Di Costanzo Pereira, Catalina Gallardo e Ignacio Yenarópulos
\end{center}

En Israel las escuelas siguen la ``regla de Maimónides'': las clases no deben tener más de 40 alumnos, así que cuando la inscripción de una cohorte supera 40 se divide en dos clases. Usamos datos de 938 clases de 5º grado (con inscripciones de entre 8 y 80 alumnos) para estimar el efecto del tamaño de la clase sobre las notas promedio de matemática (media 66, desvío 10) y de lengua (media 74, desvío 8.4). Los datos están a nivel de clase y las escuelas con más de 40 inscriptos aportan dos clases. Como la base no tiene un identificador de escuela y el porcentaje de alumnos desfavorecidos (\texttt{tip\_a}) es una variable de la escuela, aproximamos cada escuela con las combinaciones únicas de inscripción y \texttt{tip\_a} (542 escuelas) y clusterizamos los errores estándar a ese nivel.

\section*{(a) y (b) Regresiones por MCO}

\begin{table}[H]\centering
\caption{Tamaño de la clase y notas (MCO)}\label{tab:ols}
\tabla{"../output/Tabla 1.tex"}
\notas{Variables dependientes: nota promedio de la clase en matemática y en lengua. Errores estándar clusterizados por escuela entre paréntesis. * $p<0.10$, ** $p<0.05$, *** $p<0.01$.}
\end{table}

\textbf{(a)} Sin controles, el coeficiente del tamaño de la clase es positivo (columnas 1 y 2 de la Tabla~\ref{tab:ols}): un alumno más por clase se asocia con 0.16 puntos más en matemática (significativo al 1\%) y 0.08 puntos más en lengua (no significativo con errores clusterizados). Las magnitudes son chicas: pasar de una clase de 20 a una de 40 alumnos se asociaría con unos 3 puntos más en matemática, alrededor de un tercio de desvío estándar. El signo es contraintuitivo, y es difícil que refleje un efecto causal. No se cumple el supuesto de exogeneidad ($E(u\mid \text{classize})=0$): el tamaño de la clase está correlacionado con otros determinantes de las notas. Las clases grandes suelen estar en escuelas grandes y urbanas, con alumnos de mayor nivel socioeconómico, mientras que las escuelas con alumnos más desfavorecidos tienden a tener clases más chicas, en parte por políticas compensatorias. Ambas cosas sesgan el coeficiente hacia arriba.

\textbf{(b)} Al controlar por el porcentaje de alumnos desfavorecidos y por la inscripción, el coeficiente cambia de signo y deja de ser significativo ($-0.05$ en matemática y $-0.06$ en lengua; columnas 3 y 4). El porcentaje de desfavorecidos tiene un efecto negativo, grande y muy significativo sobre las notas y está correlacionado negativamente con el tamaño de la clase, lo que confirma que el coeficiente de (a) estaba sesgado por variables omitidas. De todas formas, aun con controles es difícil sostener que el tamaño de la clase sea exógeno, porque puede haber otras variables omitidas.

\section*{(c) Diseño de regresión discontinua}

La regla de Maimónides genera variación en el tamaño de la clase que, alrededor del umbral, es como si fuera aleatoria. Una escuela con 40 inscriptos tiene una clase de 40 alumnos, mientras que una con 41 tiene dos clases de alrededor de 20. Si las escuelas con inscripciones apenas por debajo y apenas por encima de 40 son comparables en todo lo demás, la diferencia en sus notas se puede atribuir al salto en el tamaño de la clase. La variable de asignación es la inscripción ($\text{enroll}$) y el umbral está entre 40 y 41: tratamos como ``por encima del umbral'' a las escuelas con 41 o más inscriptos ($c=40.5$).

El diseño es \textbf{fuzzy} y no sharp. La variable de interés es el tamaño de la clase, que no es una función determinística de cruzar el umbral: la regla cambia de forma discontinua la probabilidad (o el nivel esperado) del tratamiento, pero en general hay escuelas que dividen clases antes de llegar a 40 o que no las dividen al superar el umbral. Además, el tratamiento es continuo, de modo que lo que identifica el RD es el cociente entre el salto en las notas y el salto en el tamaño de la clase. En nuestros datos el cumplimiento de la regla es alto: el tamaño de la clase es igual a la inscripción hasta 40 y cae a cerca de la mitad a partir de 41 (Figura~\ref{fig:fs}).

Con $D_i=\mathbb{1}\{\text{enroll}_i\ge41\}$ y $\tilde x_i=\text{enroll}_i-40.5$, las especificaciones relevantes son:
\begin{align*}
\text{Primera etapa:}\quad & \text{classize}_i=\alpha_0+\alpha_1 D_i+f_0(\tilde x_i)+D_i\,f_1(\tilde x_i)+\nu_i\\
\text{Forma reducida:}\quad & y_i=\gamma_0+\gamma_1 D_i+g_0(\tilde x_i)+D_i\,g_1(\tilde x_i)+\eta_i\\
\text{RD fuzzy:}\quad & \tau=\frac{\lim_{x\downarrow c}E[y\mid x]-\lim_{x\uparrow c}E[y\mid x]}{\lim_{x\downarrow c}E[\text{classize}\mid x]-\lim_{x\uparrow c}E[\text{classize}\mid x]}=\frac{\gamma_1}{\alpha_1},
\end{align*}
donde $f$ y $g$ son polinomios de grado $p$ estimados localmente, dentro de una ventana $|\tilde x_i|\le h$ y con kernel triangular. El estimador de $\tau$ equivale a estimar por MCO en dos etapas, dentro de la ventana, $y_i=\beta_0+\tau\,\text{classize}_i+g_0(\tilde x_i)+D_i\,g_1(\tilde x_i)+\varepsilon_i$ usando $D_i$ como instrumento de $\text{classize}_i$. $\tau$ es un efecto local para las escuelas con inscripciones cercanas a 40.

\section*{(d) Pruebas de validez}

Hacemos dos pruebas estándar. La primera es un \textbf{test de densidad} de la variable de asignación \citep{cjm}: si las escuelas pudieran manipular su inscripción para quedar de un lado del umbral, veríamos un salto en la cantidad de escuelas justo antes o justo después de 40.5. La segunda es un \textbf{test de placebo con covariables predeterminadas}: las características que no pueden verse afectadas por el tratamiento, como el porcentaje de alumnos desfavorecidos, deberían ser continuas en el umbral.

\begin{figure}[H]\centering
\begin{minipage}{0.49\linewidth}\centering
\figura[\linewidth]{../output/RD_densidad.png}\\(a) Densidad de la inscripción (una observación por escuela)
\end{minipage}\hfill
\begin{minipage}{0.49\linewidth}\centering
\figura[\linewidth]{../output/RD_covariable.png}\\(b) Porcentaje de alumnos desfavorecidos
\end{minipage}
\caption{Pruebas de validez del diseño}\label{fig:valid}
\end{figure}

El test de densidad hay que hacerlo con una observación por escuela. Con los datos a nivel de clase la densidad salta mecánicamente en el umbral, porque cada escuela con más de 40 inscriptos aporta dos clases: a nivel de clase el test rechaza ($T=2.88$, p-valor $=0.004$), pero eso no indica manipulación. A nivel de escuela, el gráfico de la Figura~\ref{fig:valid}(a) muestra más escuelas justo por encima de 40 que justo por debajo, pero el test formal no rechaza la continuidad de la densidad ($T=0.60$, p-valor $=0.55$). Hay que tener en cuenta que hay pocas escuelas cerca del umbral, así que el test tiene poco poder. En el porcentaje de alumnos desfavorecidos se ve un salto hacia arriba de alrededor de 5 puntos porcentuales (Figura~\ref{fig:valid}(b)), pero no es estadísticamente significativo. En conjunto, las pruebas no rechazan la validez del diseño, aunque el patrón de los gráficos pide cautela.

¿Podría haber manipulación de la inscripción? Hay razones para pensar que sí sería posible: los directores podrían inscribir (o reportar) algunos alumnos más para pasar el umbral y conseguir una clase y un docente adicionales, y los padres más informados podrían elegir escuelas con inscripciones justo por encima de 40 para que sus hijos estén en clases más chicas. Si esto ocurriera, las escuelas justo por encima del umbral serían sistemáticamente distintas (por ejemplo, con directores más activos o familias más involucradas) y el RD estaría sesgado. Sin embargo, la inscripción depende del tamaño de la cohorte de niños del área, que es difícil de controlar con precisión, y los tests formales no detectan ni acumulación de escuelas justo por encima del umbral ni saltos en las características de los alumnos. Aun así, los gráficos muestran algo más de escuelas y algo más de alumnos desfavorecidos justo por encima de 40, y \citet{angrist2019} documentan manipulación de la inscripción alrededor de los umbrales en datos israelíes más recientes. Por eso no podemos descartar del todo la manipulación, aunque en estos datos la evidencia no es concluyente.

\section*{(e) Estimación por RD (continuity-based approach)}

\begin{figure}[H]\centering
\figura[0.55\linewidth]{../output/RD_primera_etapa.png}
\caption{Primera etapa: tamaño de la clase según la inscripción}\label{fig:fs}
\end{figure}

La primera etapa es muy fuerte: en el umbral el tamaño de la clase cae en unos 20 alumnos, de alrededor de 40 a alrededor de 20 (Figura~\ref{fig:fs}). La Tabla~\ref{tab:rd} presenta las estimaciones del RD fuzzy con polinomios locales de grado 1 a 4, con el bandwidth óptimo (que minimiza el error cuadrático medio), kernel triangular, \texttt{tip\_a} como covariable y errores clusterizados por escuela. Las Figuras~\ref{fig:rdm} y~\ref{fig:rdv} muestran los gráficos de RD (forma reducida) con ajustes polinomiales de grado 1 a 4.

\begin{table}[H]\centering
\caption{Efecto del tamaño de la clase: RD fuzzy con polinomios de grado 1 a 4}\label{tab:rd}
\scriptsize
\tabla{"../output/Tabla 2.tex"}
\notas{Cada columna es una estimación de \texttt{rdrobust} con \texttt{fuzzy(classize)}, umbral en 40.5, bandwidth óptimo MSE y kernel triangular, controlando por \texttt{tip\_a}. Coeficiente y error estándar convencionales (clusterizados por escuela); se reporta el p-valor robusto con corrección de sesgo. * $p<0.10$, ** $p<0.05$, *** $p<0.01$ (convencional).}
\end{table}

\begin{figure}[H]\centering
\figura[0.85\linewidth]{../output/RD_y1_polinomios.png}
\caption{Gráficos de RD para la nota de matemática, polinomios de grado 1 a 4}\label{fig:rdm}
\end{figure}

\begin{figure}[H]\centering
\figura[0.85\linewidth]{../output/RD_y2_polinomios.png}
\caption{Gráficos de RD para la nota de lengua, polinomios de grado 1 a 4}\label{fig:rdv}
\end{figure}

Con el polinomio lineal, el efecto sobre la nota de matemática es prácticamente nulo ($-0.09$ puntos por alumno, con un p-valor robusto de 0.76), mientras que en lengua un alumno más por clase reduce la nota en 0.28 puntos (p-valor robusto de 0.036). Como en el umbral la clase pasa de unos 40 a unos 20 alumnos, el efecto implícito de la reducción es de alrededor de 5.6 puntos en lengua, unos dos tercios de desvío estándar. Con polinomios de mayor grado las estimaciones puntuales son similares o algo mayores (en lengua, entre $-0.31$ y $-0.43$), pero los errores estándar aumentan y la significatividad cae: con polinomio de grado 2 el efecto es significativo solo al 10\% y con grados 3 y 4 deja de serlo según los p-valores robustos. En matemática el efecto nunca es significativo.

Preferimos el polinomio de grado 1 (regresión lineal local). Los polinomios de grado alto tienden a sobreajustar los datos, son muy sensibles a las observaciones cercanas al borde del soporte (justamente el umbral) y generan estimaciones ruidosas y poco confiables \citep{gelman}. El ajuste lineal local con bandwidth óptimo es el estándar recomendado por \citet{cattaneo}. Los gráficos lo ilustran: los ajustes de grado 3 y 4 son muy irregulares, con curvas que caen bruscamente justo antes del umbral para seguir a unas pocas observaciones.

\section*{(f) Distintos bandwidths}

\begin{table}[H]\centering
\caption{Efecto del tamaño de la clase: RD fuzzy lineal con bandwidth 7, 10 y 15}\label{tab:bw}
\scriptsize
\tabla{"../output/Tabla 3.tex"}
\notas{Ver notas de la Tabla~\ref{tab:rd}. Polinomio local de grado 1 con el bandwidth indicado.}
\end{table}

\begin{figure}[H]\centering
\begin{minipage}{0.32\linewidth}\centering\figura[\linewidth]{../output/RD_y1_h7.png}\\ Matemática, $h=7$\end{minipage}\hfill
\begin{minipage}{0.32\linewidth}\centering\figura[\linewidth]{../output/RD_y1_h10.png}\\ Matemática, $h=10$\end{minipage}\hfill
\begin{minipage}{0.32\linewidth}\centering\figura[\linewidth]{../output/RD_y1_h15.png}\\ Matemática, $h=15$\end{minipage}\\[0.5em]
\begin{minipage}{0.32\linewidth}\centering\figura[\linewidth]{../output/RD_y2_h7.png}\\ Lengua, $h=7$\end{minipage}\hfill
\begin{minipage}{0.32\linewidth}\centering\figura[\linewidth]{../output/RD_y2_h10.png}\\ Lengua, $h=10$\end{minipage}\hfill
\begin{minipage}{0.32\linewidth}\centering\figura[\linewidth]{../output/RD_y2_h15.png}\\ Lengua, $h=15$\end{minipage}
\caption{Gráficos de RD con bandwidth 7, 10 y 15}\label{fig:bw}
\end{figure}

La elección del bandwidth implica un trade-off entre sesgo y varianza. Con un bandwidth chico se comparan solo escuelas muy parecidas, cercanas al umbral, lo que reduce el sesgo de aproximar la relación con un polinomio lineal, pero se usan pocas observaciones y la estimación es menos precisa. Con un bandwidth grande se gana precisión, pero se incluyen escuelas lejos del umbral y, si la relación no es lineal en todo el rango, aumenta el sesgo. Esperamos entonces estimaciones parecidas y errores estándar que caen a medida que aumenta el bandwidth. Eso es lo que encontramos (Tabla~\ref{tab:bw}). En matemática el efecto es chico y no significativo con los tres bandwidths ($-0.09$, $-0.12$ y $-0.17$). En lengua las estimaciones son muy estables ($-0.30$, $-0.28$ y $-0.26$) y la precisión mejora con el bandwidth: el p-valor robusto pasa de 0.10 con $h=7$ y 0.08 con $h=10$ a 0.035 con $h=15$. El bandwidth óptimo (alrededor de 8 a 10) queda en el medio de este rango.

\section*{(g) La regla de Maimónides como variable instrumental}

\begin{table}[H]\centering
\caption{Tamaño de la clase y notas: IV con la regla de Maimónides}\label{tab:iv}
\tabla{"../output/Tabla 4.tex"}
\notas{MC2E: \textit{classize} instrumentado con $z_s=e_s/(\lfloor(e_s-1)/40\rfloor+1)$. Todas las columnas controlan por \texttt{tip\_a}. Errores estándar clusterizados por escuela entre paréntesis. * $p<0.10$, ** $p<0.05$, *** $p<0.01$.}
\end{table}

El tamaño de clase predicho por la regla está muy correlacionado con el tamaño efectivo (correlación de 0.97), así que el instrumento es muy relevante. Sin controlar por la inscripción (columnas 1 y 3 de la Tabla~\ref{tab:iv}), el efecto en matemática es chico y no significativo ($-0.02$) y en lengua es de $-0.11$ puntos por alumno, significativo al 1\%. Pero $z_s$ es una función de la inscripción, de modo que sin controlar por ella el instrumento mezcla el salto que genera la regla con la relación suave entre el tamaño de la escuela y las notas. Por eso \citet{angrist} controlan también por la inscripción. Con ese control (columnas 2 y 4), ambos efectos son negativos ($-0.05$ en matemática y $-0.07$ en lengua) pero no significativos. Las estimaciones de IV son menores en valor absoluto que las del RD fuzzy porque usan todo el rango de inscripciones con una tendencia lineal, en lugar de concentrarse alrededor del umbral. Además, en esta muestra hay un solo umbral (inscripciones de hasta 80), mientras que \citet{angrist} explotan varios umbrales (80, 120, etc.), lo que les da más variación y más precisión. En conjunto, la evidencia sugiere que reducir el tamaño de la clase mejora las notas de lengua, mientras que para matemática no encontramos efectos significativos.

\begin{thebibliography}{9}
\bibitem[Angrist y Lavy(1999)]{angrist} Angrist, J. D., y Lavy, V. (1999). Using Maimonides' rule to estimate the effect of class size on scholastic achievement. \textit{Quarterly Journal of Economics}, 114(2), 533--575.
\bibitem[Angrist et al.(2019)]{angrist2019} Angrist, J. D., Lavy, V., Leder-Luis, J., y Shany, A. (2019). Maimonides' rule redux. \textit{American Economic Review: Insights}, 1(3), 309--324.
\bibitem[Cattaneo, Idrobo y Titiunik(2020)]{cattaneo} Cattaneo, M. D., Idrobo, N., y Titiunik, R. (2020). \textit{A Practical Introduction to Regression Discontinuity Designs: Foundations}. Cambridge University Press.
\bibitem[Cattaneo, Jansson y Ma(2020)]{cjm} Cattaneo, M. D., Jansson, M., y Ma, X. (2020). Simple local polynomial density estimators. \textit{Journal of the American Statistical Association}, 115(531), 1449--1455.
\bibitem[Gelman e Imbens(2019)]{gelman} Gelman, A., e Imbens, G. (2019). Why high-order polynomials should not be used in regression discontinuity designs. \textit{Journal of Business \& Economic Statistics}, 37(3), 447--456.
\end{thebibliography}

\newpage
\appendix
\section*{Apéndice: código}
\lstinputlisting{../programs/ps10.do}

\end{document}
